\documentclass[preprint,12pt]{elsarticle}

\journal{Computer Methods and Programs in Biomedicine}
\raggedbottom

\usepackage[T1]{fontenc}
\usepackage{amsmath,amssymb}
\usepackage{booktabs}
\usepackage{array}
\usepackage{longtable}
\usepackage{pdflscape}
\usepackage{microtype}
\usepackage{caption}
\usepackage{algorithm}
\usepackage{algpseudocode}
\usepackage{placeins}
\usepackage{url}

\captionsetup[table]{skip=4pt}
\captionsetup[longtable]{skip=4pt}

\setlength{\textfloatsep}{8pt plus 2pt minus 2pt}
\setlength{\floatsep}{8pt plus 2pt minus 2pt}
\setlength{\intextsep}{8pt plus 2pt minus 2pt}

\newcolumntype{L}[1]{>{\raggedright\arraybackslash}p{#1}}

\renewcommand{\thetable}{S\arabic{table}}
\renewcommand{\thealgorithm}{S\arabic{algorithm}}

\begin{document}

\begin{frontmatter}
\title{Supplementary File 2. Supplementary Methods}
\end{frontmatter}

\subsection*{S1. Metadata-governance admissibility layer}

MIABIS supports interoperable metadata representation and specimen-level contextualization, ISO 20387 and ISBER Best Practices support governance, quality management, documentation, monitoring, and traceability requirements, and SPREC preserves the structured pre-analytical trajectory of fluid and solid biospecimens. Therefore, the governance layer evaluates whether a specimen is admissible for scoring, whereas the SPREC-derived layer supports the subsequent computation of the RCS. The metadata-governance admissibility layer is summarized in Supplementary Table~S1.

\begin{landscape}
\scriptsize
\setlength{\tabcolsep}{3pt}
\renewcommand{\arraystretch}{1.10}

\begin{longtable}{@{}L{4.0cm}L{10.0cm}L{7.5cm}@{}}
\caption{Metadata-governance admissibility layer.}
\label{tab:governance_layers}\\

\toprule
\textbf{Evidence category} & \textbf{Model input} & \textbf{Model effect} \\
\midrule
\endfirsthead
\toprule
\textbf{Evidence category} & \textbf{Model input} & \textbf{Model effect} \\
\midrule
\endhead

\midrule
\multicolumn{3}{r}{\textit{Continued on next page}}\\
\endfoot

\bottomrule
\endlastfoot

\multicolumn{3}{@{}l}{\textbf{MIABIS-derived metadata layer}} \\
\midrule

Sample-level metadata
& Sample identifier, sample type, matrix type, and local sample record
& Required for specimen identification and matrix assignment \\

Donor or source linkage
& Link between the sample and its donor, source entity, or originating subject record
& Required for contextual interpretation and traceability \\

Collection context
& Biobank, collection, or research-resource context associated with the specimen
& Required for interoperable specimen qualification \\

Event-level metadata
& Collection, processing, storage, shipment, or other time- and place-dependent events
& Required for reconstructing the specimen trajectory \\

Metadata completeness
& Presence of mandatory descriptive fields needed for interpretation and sharing
& Missing critical metadata contributes to governance-gate failure \\

\midrule
\multicolumn{3}{@{}l}{\textbf{ISO 20387-derived quality and traceability layer}} \\
\midrule

Quality management
& Documented quality procedures, operational consistency, quality control, and competence-related records
& Supports admissibility and auditability \\

Traceability
& Unique identification, chain-of-custody evidence, location history, and process documentation
& Required for governance-gate approval \\

Nonconformity control
& Records of deviations, nonconformities, corrective actions, or technical restrictions
& Critical nonconformities trigger review or restriction \\

Data reliability
& Consistency between sample record, processing history, storage condition, and associated metadata
& Required for reliable downstream qualification \\

\midrule
\multicolumn{3}{@{}l}{\textbf{ISBER-informed operational governance layer}} \\
\midrule

Biospecimen handling
& Documented collection, processing, storage, retrieval, and distribution practices
& Supports operational admissibility and reproducibility \\

Storage monitoring
& Temperature monitoring, cold-chain evidence, freezer status, and storage-condition records
& Required when thermal control is relevant \\

Documentation and SOPs
& Standard operating procedures, processing records, and documented handling conditions
& Supports governance-gate approval and technical review \\

Deviation documentation
& Recorded handling deviations, storage failures, or processing exceptions
& Critical deviations may trigger review or quarantine \\

\midrule
\multicolumn{3}{@{}l}{\textbf{Controlled terminology and context validation}} \\
\midrule

Semantic validity and context compatibility
& Valid SPREC or controlled terminology values compatible with matrix type, event type, and recorded specimen context
& Invalid, unrecognized, or severely incompatible codes contribute to governance-gate failure or technical review \\

\addlinespace[0.4em]
\multicolumn{3}{@{}p{20.5cm}@{}}{\textit{Table note.} MIABIS-derived metadata, ISO 20387 requirements, ISBER Best Practices, and terminology governance do not 
define operational penalties; instead, they determine whether the specimen is admissible for SPREC-based RCS computation.} \\

\end{longtable}
\normalsize
\end{landscape}

Context-validation procedures were included because the RCS depends on SPREC-coded axes, while a specimen can only be scored when its coded values 
are valid, recognized, and compatible with the recorded matrix type and event context. Invalid or severely incompatible codes were treated as 
governance-admissibility failures rather than as compensable operational penalties. Accordingly, values classified as not eligible for additive 
scoring are not interpreted as zero-risk conditions; they indicate insufficient semantic, documentary, or contextual validity for score computation.

\subsection*{S2. Axis-level penalty structure and biological-operational rationale}

In practical terms, MIABIS-derived metadata were used to represent the minimum descriptive and interoperability context required for specimen-level 
interpretation. This includes sample identifiers, sample type, donor or source linkage, collection or biobank context, and event-level information 
related to collection, processing, storage, shipment, or other time- and place-dependent operations. ISO 20387 and ISBER Best Practices informed the 
governance dimension of the model, including quality management, documentation, traceability, monitoring, deviation records, nonconformity control, 
and operational consistency.

Therefore, the final qualification framework generated by the model is not a native MIABIS class, nor a direct output of SPREC, ISO 20387, or ISBER 
Best Practices. Instead, it is a controlled local qualification layer generated by the semantic validation layer and represented alongside the MIABIS-
compatible metadata, traceability descriptors, documentation records, monitoring evidence, and terminology-validation elements summarized in 
Supplementary Table~S1.

Each observed SPREC code then received a severity value within its own domain. Thus, the same SPREC letter may have different meanings across 
different axes and was interpreted only within its corresponding pre-analytical context. This approach preserves SPREC as a standardized biological 
descriptive structure while allowing the RCS to quantify the accumulated operational burden of governance-admissible specimens. The priority of each 
SPREC-derived axis, represented by \(W_i^{(k)}\), was defined from the operational rationale summarized in Supplementary Table~S2. 

\begin{landscape}
\scriptsize
\setlength{\tabcolsep}{3.5pt}
\renewcommand{\arraystretch}{1.12}

\begin{longtable}{@{}L{3.7cm}L{9.2cm}L{1.3cm}L{7.4cm}@{}}

\caption{Axis-level penalty structure for fluid and solid biospecimens.}
\label{tab:axis_penalty_structure}\\[-0.3em]

\toprule
\textbf{SPREC-derived axis}
& \textbf{Biological-operational rationale}
& \textbf{Weight}
& \textbf{Severity interpretation} \\
\midrule
\endfirsthead

\toprule
\textbf{SPREC-derived axis}
& \textbf{Biological-operational rationale}
& \textbf{Weight}
& \textbf{Severity interpretation} \\
\midrule
\endhead

\midrule
\multicolumn{4}{r}{\textit{Continued on next page}}\\
\endfoot

\bottomrule
\endlastfoot

\multicolumn{4}{@{}l}{\textbf{Fluid biospecimens}} \\
\midrule

Pre-centrifugation delay
& May affect cellular metabolism, hemolysis-related artifacts, and molecular stability before separation \cite{moore2011brisq,hassis2015,lehmann2012}.
& 30
& Severity increases with longer admissible delay intervals. \\

Primary centrifugation
& May influence plasma or serum composition through incomplete separation or cellular contamination \cite{hassis2015,lehmann2012}.
& 15
& Severity increases with non-standard admissible centrifugation conditions. \\

Second centrifugation
& Relevant for plasma clarification and reduction of residual cellular contamination \cite{hassis2015,lehmann2012}.
& 10
& Severity increases when workflow-compatible clarification is absent or suboptimal. \\

Post-centrifugation delay
& May influence ex vivo molecular stability after sample separation \cite{moore2011brisq,hassis2015}.
& 20
& Severity increases with longer admissible post-processing delay. \\

Long-term storage
& Storage conditions influence long-term molecular preservation and analytical reproducibility \cite{moore2011brisq,hassis2015,lehmann2012}.
& 25
& Severity increases with less stable admissible storage conditions. \\

\addlinespace[0.4em]
\midrule
\multicolumn{4}{@{}l}{\textbf{Solid biospecimens}} \\
\midrule

Warm ischemia time
& May induce tissue stress responses and early molecular degradation before preservation \cite{moore2011brisq,huang2001}.
& 25
& Severity increases with longer admissible warm ischemia intervals. \\

Cold ischemia time
& May affect RNA integrity and tissue-specific molecular preservation before stabilization \cite{huang2001,viana2013}.
& 25
& Severity increases with longer admissible cold ischemia intervals. \\

Fixation or stabilization type
& Preservation method affects compatibility with downstream analytical workflows \cite{moore2011brisq,lehmann2012}.
& 15
& Severity increases with less compatible admissible preservation methods. \\

Fixation time
& Fixation duration may affect morphology, antigen preservation, and nucleic-acid recovery \cite{moore2011brisq,lehmann2012}.
& 20
& Severity increases with greater deviation from the reference fixation interval. \\

Long-term storage
& Storage conditions influence long-term tissue preservation and reproducibility \cite{moore2011brisq,lehmann2012}.
& 15
& Severity increases with less stable admissible storage conditions. \\

\addlinespace[0.5em]
\multicolumn{4}{@{}p{21.6cm}@{}}{
\textit{Table note.}
The effective penalty depends on the intra-axis severity function \(s_i(x_i)\), constrained between 0 and 1 and detailed in Supplementary File~1, Supplementary Table~S3. Unknown, missing, invalid, or semantically incompatible values are evaluated through the non-compensable governance gate.
} \\

\end{longtable}

\normalsize
\end{landscape}


In contrast, the intra-axis severity value \(s_i(x_i)\) was derived from the specific condition represented by each SPREC code. For temporal SPREC axes, severity followed the time intervals defined in SPREC 2.0. For non-temporal axes, such as centrifugation, fixation or stabilization type, and long-term storage, severity followed workflow compatibility, preservation stability, and relevance of the coded condition. The complete intra-axis severity mapping is provided in Supplementary File~1, Supplementary Table~S3.

A coded condition received low severity when it was compatible with the intended workflow but represented a less stable or less optimal processing or 
preservation state. Conversely, the same coded condition received maximum severity when it was incompatible with the intended analytical workflow or 
when an expected processing step was absent. Therefore, Supplementary Methods Table~S3 does not repeat the weighting rationale; instead, it summarizes how the weighted SPREC-derived variables enter the fluid- and solid-specific RCS equations.

\begin{landscape}
\scriptsize
\setlength{\tabcolsep}{2pt}
\renewcommand{\arraystretch}{1.05}

\setcounter{table}{2}

\begin{longtable}{@{}L{3.0cm}L{2.0cm}L{4.8cm}L{6.0cm}@{}}
\caption{Matrix-specific RCS calculation variables.}
\label{tab:rcs_calculation_variables}\\

\toprule
\textbf{SPREC-derived axis} 
& \textbf{RCS variable} 
& \textbf{Penalty term} 
& \textbf{Role in RCS computation} \\
\midrule
\endfirsthead

\toprule
\textbf{SPREC-derived axis} 
& \textbf{RCS variable} 
& \textbf{Penalty term} 
& \textbf{Role in RCS computation} \\
\midrule
\endhead

\midrule
\multicolumn{4}{r}{\textit{Continued on next page}}\\
\endfoot

\bottomrule
\endlastfoot

\multicolumn{4}{@{}l}{\textbf{Fluid biospecimens}} \\
\midrule

Pre-centrifugation delay
& \(P_{pre}\)
& \(W_{pre}^{(fluid)} s_{pre}(x_{pre})\)
& Weighted contribution of delay before centrifugation \\

Primary centrifugation
& \(P_{cent1}\)
& \(W_{cent1}^{(fluid)} s_{cent1}(x_{cent1})\)
& Weighted contribution of the primary centrifugation condition \\

Second centrifugation
& \(P_{cent2}\)
& \(W_{cent2}^{(fluid)} s_{cent2}(x_{cent2})\)
& Weighted contribution of second centrifugation when applicable \\

Post-centrifugation delay
& \(P_{post}\)
& \(W_{post}^{(fluid)} s_{post}(x_{post})\)
& Weighted contribution of delay after centrifugation \\

Long-term storage
& \(P_{store}\)
& \(W_{store}^{(fluid)} s_{store}(x_{store})\)
& Weighted contribution of storage stability \\

\midrule
\multicolumn{4}{@{}l}{\textbf{Fluid-specific accumulated penalty}} \\
\midrule

\multicolumn{2}{@{}L{5.0cm}}{\(P_{bio}^{(fluid)}\)}
& \multicolumn{2}{L{10.8cm}@{}}{
\(P_{pre}+P_{cent1}+P_{cent2}+P_{post}+P_{store}\)
} \\

\midrule
\multicolumn{4}{@{}l}{\textbf{Solid biospecimens}} \\
\midrule

Warm ischemia time
& \(P_{warm}\)
& \(W_{warm}^{(solid)} s_{warm}(x_{warm})\)
& Weighted contribution of warm ischemia exposure \\

Cold ischemia time
& \(P_{cold}\)
& \(W_{cold}^{(solid)} s_{cold}(x_{cold})\)
& Weighted contribution of cold ischemia exposure \\

Fixation or stabilization type
& \(P_{fix}\)
& \(W_{fix}^{(solid)} s_{fix}(x_{fix})\)
& Weighted contribution of preservation-method compatibility \\

Fixation time
& \(P_{fixTime}\)
& \(W_{fixTime}^{(solid)} s_{fixTime}(x_{fixTime})\)
& Weighted contribution of fixation-duration deviation \\

Long-term storage
& \(P_{store}\)
& \(W_{store}^{(solid)} s_{store}(x_{store})\)
& Weighted contribution of storage stability \\

\midrule
\multicolumn{4}{@{}l}{\textbf{Solid-specific accumulated penalty}} \\
\midrule

\multicolumn{2}{@{}L{5.0cm}}{\(P_{bio}^{(solid)}\)}
& \multicolumn{2}{L{10.8cm}@{}}{
\(P_{warm}+P_{cold}+P_{fix}+P_{fixTime}+P_{store}\)
} \\

\addlinespace[0.4em]
\multicolumn{4}{@{}L{15.8cm}@{}}{
\textit{Table note.}
Each variable represents the weighted penalty contribution of one SPREC-derived axis after governance admissibility has been verified. This table should be read together with Supplementary Methods Table~S2. The term \(W_i^{(k)}\) represents the maximum contribution of axis \(i\) for matrix \(k\), whereas \(s_i(x_i)\) represents the normalized intra-axis severity of the observed SPREC-coded condition. Unknown, missing, invalid, or semantically incompatible values are not represented as additive penalty terms and are evaluated through the non-compensable governance gate.
} \\

\end{longtable}

\normalsize
\end{landscape}

For temporal derived axes, values of \(s_i(x_i)\) were assigned according to the time intervals defined in SPREC 2.0, whereas non-temporal axes were assigned according to workflow compatibility, preservation stability, and operational relevance of the coded condition. Unknown, invalid, missing, or semantically incompatible values were not treated as additive penalties; instead, they were evaluated through the non-compensable governance gate.

Accordingly, \(W_i^{(k)}\) defines the maximum contribution of each axis to \(P_{bio}^{(k)}\), whereas \(s_i(x_i)\) defines the normalized severity of the observed coded condition within that axis in SPREC. Values marked as ``Not scored'' were not interpreted as zero-risk or optimal conditions. Instead, they were excluded from additive severity assignment because SPREC is a descriptive pre-analytical coding system rather than a numerical quality score. Unresolved coded values were therefore treated as failures of semantic admissibility, traceability, documentation, or contextual validity rather than as compensable operational deviations \cite{lehmann2012,merinomartinez2016,campbell2018,iso20387}.

For workflow-dependent or conditional SPREC-derived conditions, severity was assigned only when the intended downstream workflow and compatibility context were documented. This rule was particularly relevant for centrifugation and storage conditions, because plasma and cell-free workflows require adequate removal of cellular contaminants, whereas long-term storage suitability depends on the biomolecule, matrix, temperature range, preservation method, and intended downstream application \cite{hassis2015,binderup2016,mendy2017iarc}. When this context was absent, the condition was routed to the non-compensable governance gate.

The priority of each derived axis was represented by the maximum contribution \(W_i^{(k)}\). The sum of the penalty-generating axes was constrained to 100 points within each matrix, preserving the interpretability of the RCS as a 0--100 score. The intra-axis severity value \(s_i(x_i)\) converts each observed coded condition into a normalized value between 0 and 1.

The deterministic implementation of the RCS framework is described in Supplementary Methods Algorithm~S1, which presents the governance-aware certification workflow applied to a single biospecimen profile. The parallel implementation, described in Supplementary Methods Algorithm~S2, preserves the deterministic semantics of the RCS framework while allowing independent biospecimen profiles to be distributed across worker threads. Algorithm~S3 presents the generic axis-level penalty function used by the multithread-capable scoring engine.

\subsection*{S3. Deterministic and multithreaded RCS certification algorithms}

The procedure begins with governance admissibility assessment, followed by matrix-specific axis selection, SPREC-derived intra-axis severity resolution, weighted biological-operational penalty accumulation, and deterministic A--E certification assignment.

\begin{algorithm}[!htbp]
\caption{Governance-aware Ribeiro Classification Score certification algorithm}
\label{alg:rcs_certification}
\footnotesize
\begin{algorithmic}[1]

\Procedure{CertifyBiospecimen}{$\mathrm{matrixType}, sprecCodes, metadata, telemetry$}
\State $S_{max} \gets 100$
\State $G_{gov} \gets$ \Call{EvaluateGovernanceGate}{$sprecCodes, metadata, telemetry$}

\If{$G_{gov} = 0$}
    \State \Return Grade E, Governance-failure condition
\EndIf

\State $axisSet \gets$ \Call{SelectMatrixSpecificAxes}{$\mathrm{matrixType}$}
\State $P_{bio} \gets 0$

\ForAll{$axis_i$ in $axisSet$}
    \State $W_i \gets$ \Call{MaximumAxisWeight}{$axis_i, \mathrm{matrixType}$}
    \State $x_i \gets$ \Call{ObservedSPRECCode}{$axis_i, sprecCodes$}
    \State $s_i \gets$ \Call{DomainSpecificSeverity}{$axis_i, x_i$}
    \State $P_{bio} \gets P_{bio} + W_i \cdot s_i$
\EndFor

\State $RCS \gets S_{max} - P_{bio}$

\If{$RCS \geq 90$}
    \State \Return $RCS$, Grade A, Optimal profile
\ElsIf{$RCS \geq 80$}
    \State \Return $RCS$, Grade B, Mild suboptimal profile
\ElsIf{$RCS \geq 65$}
    \State \Return $RCS$, Grade C, Moderate suboptimal profile
\ElsIf{$RCS \geq 50$}
    \State \Return $RCS$, Grade D, Severe suboptimal profile
\Else
    \State \Return $RCS$, Grade E, Critical penalty-burden condition
\EndIf

\EndProcedure
\end{algorithmic}
\normalsize
\end{algorithm}

Because each biospecimen profile is evaluated independently after governance admissibility assessment, the computational architecture supports scalable multithreaded execution for batch-oriented processing \cite{dean2008mapreduce,openmp2026}.

Algorithm~\ref{alg:parallel_rcs_certification} illustrates a hierarchical parallel implementation in which independent biospecimen profiles are distributed across worker threads while preserving deterministic scoring semantics, governance-aware admissibility logic, and reproducible certification behavior. Algorithm~\ref{alg:axis_penalty_function} presents the generic axis-level penalty function used by the multithread-capable scoring engine.

\begin{algorithm}[!htbp]
\caption{Hierarchical multithreaded governance-aware RCS certification}
\label{alg:parallel_rcs_certification}
\footnotesize
\begin{algorithmic}[1]

\Procedure{ParallelCertifyBatch}{$biospecimenBatch, nThreads$}

\State Partition $biospecimenBatch$ into $nThreads$ independent subsets

\ForAll{worker thread $t$ in parallel}
    \State Initialize \Call{ThreadLocalStore}{$t$}

    \ForAll{biospecimen $b$ in subset $t$}

        \State $G_{gov} \gets$ \Call{EvaluateGovernanceGate}{$b.sprecCodes, b.metadata, b.telemetry$}

        \If{$G_{gov} = 0$}
            \State \Call{ThreadLocalStore}{$t, b.id, \mathrm{Grade~E}, \mathrm{Governance\ failure}$}
        \Else
            \If{$b.\mathrm{matrixType} = \mathrm{fluid}$}
                \State $axisSet \gets \{P_{pre}, P_{cent1}, P_{cent2}, P_{post}, P_{store}\}$
            \ElsIf{$b.\mathrm{matrixType} = \mathrm{solid}$}
                \State $axisSet \gets \{P_{warm}, P_{cold}, P_{fix}, P_{fixTime}, P_{store}\}$
            \EndIf

            \State $P_{bio}^{(k)} \gets \sum$ \Call{ParallelMap}{$\mathrm{ComputeAxisPenalty}, axisSet, b.sprecCodes, b.metadata$}
            \State $P_{bio}^{(k)} \gets \sum_{p_i^{(k)} \in \mathcal{P}^{(k)}} p_i^{(k)}$
            \State $RCS \gets 100 - P_{bio}^{(k)}$
            \State $grade \gets$ \Call{ScoreGrade}{$RCS$}
            \State \Call{ThreadLocalStore}{$t, b.id, RCS, grade, \mathrm{Score\mbox{-}based}$}
        \EndIf

    \EndFor
\EndFor

\State $classificationTable \gets$ \Call{MergeThreadLocalStores}{$1,\ldots,nThreads$}
\State \Return $classificationTable$

\EndProcedure

\end{algorithmic}

\vspace{0.5em}
\noindent\footnotesize
\textit{Algorithm note.}
\textsc{ThreadLocalStore}(\(\cdot\)) denotes a worker-specific output buffer used to store classification results locally within each thread. This avoids concurrent writes to a shared global table during parallel execution. After all worker threads complete, \textsc{MergeThreadLocalStores}(\(\cdot\)) deterministically combines the thread-local buffers into the final classification table.
\normalsize
\end{algorithm}

\begin{algorithm}[!htbp]
\caption{Axis-level RCS penalty function}
\label{alg:axis_penalty_function}
\footnotesize
\begin{algorithmic}[1]

\Procedure{ComputeAxisPenalty}{$axis_i, sprecCodes, metadata$}
    \State $W_i^{(k)} \gets$ \Call{MaximumAxisWeight}{$axis_i$}
    \State $x_i \gets$ \Call{ObservedSPRECCode}{$axis_i, sprecCodes$}
    \State $s_i(x_i) \gets$ \Call{ResolveIntraAxisSeverity}{$axis_i, x_i, metadata$}

    \If{$s_i(x_i) = \mathrm{NotScored}$}
        \State Raise governance/context review flag
    \EndIf

    \State $p_i^{(k)} \gets W_i^{(k)} \cdot s_i(x_i)$
    \State \Return $p_i^{(k)}$
\EndProcedure

\end{algorithmic}
\normalsize
\end{algorithm}

The parallel implementation does not modify the deterministic semantics of the RCS framework. Governance admissibility assessment remains non-compensable and precedes penalty computation. After governance approval, independent axis-level penalty functions may be evaluated concurrently because the weighted penalty terms are mathematically independent prior to reduction by summation.

\subsection*{S4. Synthetic validation scenarios}

The synthetic validation scenarios used for deterministic internal validation are summarized in Supplementary Table~S4.

\setcounter{table}{3}
\begin{table}[!htbp]
\caption{Synthetic validation scenarios.}
\label{tab:synthetic_validation_scenarios}
\centering

\footnotesize
\setlength{\tabcolsep}{3pt}
\renewcommand{\arraystretch}{1.08}

\begin{tabular}{@{}L{3.0cm}L{5.2cm}L{2.4cm}L{5.8cm}@{}}

\toprule
\textbf{Synthetic scenario}
& \textbf{Expected penalty or governance profile}
& \textbf{Expected grade}
& \textbf{Validation purpose} \\
\midrule

Optimal
& Minimal or null operational penalty burden under admissible governance
& Grade A
& Verify preservation of high-score profiles \\

Mild suboptimal
& Small accumulated SPREC-derived pre-analytical deviation
& Grade B
& Verify detection of limited suboptimality \\

Moderate suboptimal
& Intermediate accumulated SPREC-derived pre-analytical deviation
& Grade C
& Verify intermediate score stratification \\

Severe suboptimal
& High accumulated biological-operational penalty burden without governance failure
& Grade D
& Verify severe but still score-admissible degradation \\

Critical penalty burden
& Accumulated penalty exceeding the critical score threshold
& Grade E
& Verify Grade E assignment due to critical RCS-derived penalty burden \\

Governance failure
& Missing, invalid, inconsistent, or non-admissible governance evidence
& Grade E
& Verify Grade E assignment due to non-compensable governance failure \\

\bottomrule

\end{tabular}

\vspace{0.4em}

\begin{minipage}{0.95\textwidth}
\footnotesize
\textit{Table note.}
The scenarios were used to test whether the score function, threshold structure, and governance gate generated the expected A--E certification behavior.
\end{minipage}

\end{table}

\subsection*{S5. Computational environment and reproducibility configuration}

All validation analyses and figure-generation routines were implemented in \textsf{R} \cite{RCoreTeam2024}. The final full validation run was executed in an isolated Linux virtual-machine environment with the large-benchmark configuration enabled. The computational environment, benchmark settings, package versions, and reproducibility parameters used in the final validation workflow are summarized in Supplementary Table~S5.

\setcounter{table}{4}
\begin{table}[!htbp]
\caption{Computational environment used for the final RCS validation run.}
\label{tab:computational_environment}
\centering

\footnotesize
\setlength{\tabcolsep}{3pt}
\renewcommand{\arraystretch}{1.08}

\begin{tabular}{@{}L{4.4cm}L{11.0cm}@{}}

\toprule
\textbf{Component}
& \textbf{Final validation configuration} \\
\midrule

Execution date and time
& 2026-05-05 15:45:44 \\

Working directory
& \texttt{/home/denise/rcs\_validation/RCS\_validation\_package} \\

Operating system
& Debian GNU/Linux 12 (bookworm), x86\_64-pc-linux-gnu, 64-bit \\

R version
& R version 4.2.2 Patched (2022-11-10 r83330) \\

Benchmark mode
& \texttt{RUN\_LARGE\_BENCH=TRUE} \\

Benchmark repetitions
& \texttt{BENCH\_REPS=5} \\

Random seed
& \texttt{RCS\_SEED=20260504} \\

Main R packages attached
& \texttt{dplyr} 1.2.1;
\texttt{tidyr} 1.3.2;
\texttt{purrr} 1.2.2;
\texttt{ggplot2} 4.0.3;
\texttt{readr} 2.2.0;
\texttt{stringr} 1.6.0;
\texttt{scales} 1.4.0;
\texttt{tibble} 3.3.1;
\texttt{forcats} 1.0.1;
\texttt{broom} 1.0.12;
\texttt{patchwork} 1.3.2 \\

BLAS/LAPACK
& System BLAS \texttt{libblas.so.3.11.0}; system LAPACK \texttt{liblapack.so.3.11.0} \\

\bottomrule

\end{tabular}

\end{table}

\FloatBarrier

\bibliographystyle{elsarticle-num}
\bibliography{references}

\end{document}

